\documentclass[aspectratio=169]{beamer}
\usetheme{default}
\usecolortheme{whale}
\setbeamertemplate{navigation symbols}{}
\setbeamertemplate{footline}{}
\setbeamertemplate{headline}{}
\setbeamerfont{title}{size=\Large,series=\bfseries}
\setbeamerfont{frametitle}{size=\Large,series=\bfseries}

\usepackage{amsmath,amssymb}
\usepackage{tikz}
\usepackage{xcolor}
\usepackage{microtype}

\definecolor{ok}{HTML}{2ca02c}
\definecolor{bad}{HTML}{d62728}
\definecolor{sig}{HTML}{1f77b4}
\definecolor{grey}{HTML}{888888}

\newcommand{\bigq}[1]{\begin{center}{\Large #1}\end{center}}
\newcommand{\keybox}[1]{\begin{center}\colorbox{grey!15}{\parbox{0.92\linewidth}{\centering #1}}\end{center}}
\newcommand{\plain}[1]{\textit{\textcolor{sig}{Plain words:} #1}}

\title{The Honest Caveats, Explained}
\subtitle{From basic level to advanced, with the real numbers}
\author{ASYU 2026}
\date{}

\begin{document}

\frame{\titlepage}

\begin{frame}{How to read this deck}
\begin{itemize}
\item Assumes \textbf{only basic arithmetic}: bigger/smaller, fractions, averages.
\item Every term is defined before it is used.
\item One idea per slide. Numbers shown are the real measured values.
\item Goal: understand \emph{why} three specific caveats in the two papers
cannot be ``fixed'' by editing, because they are facts, not typos.
\end{itemize}
\vspace{1em}
\keybox{A caveat is a sentence that says ``here is exactly how far our
result does and does not reach''. Stating it is honesty, not weakness.}
\end{frame}

% =====================================================================
\section{Part 0: two kinds of problem}
% =====================================================================
\begin{frame}{Two completely different kinds of problem}
\begin{itemize}
\item \textbf{Copy-edit defect}: a number typed wrong, a figure not
regenerated, a contradictory sentence. \textcolor{ok}{Fixable by editing.}
We already fixed all of these.
\item \textbf{Empirical truth}: what the data actually shows, including
where it is weak. \textcolor{bad}{Not fixable by editing}, only by
collecting more/different evidence, or by lying.
\end{itemize}
\vspace{1em}
\keybox{The three caveats in this deck are the second kind. You change
them only by doing new science, or by being dishonest. We choose
neither: we state them plainly.}
\end{frame}

% =====================================================================
\section{Part 1: the words you need}
% =====================================================================
\begin{frame}{Precision and coverage (plain)}
\plain{Think of a gate that says ``accept'' or ``flag for a human''.}
\begin{itemize}
\item \textbf{Coverage}: of all receipts, what fraction does the gate
\emph{accept} (not send to a human). Example: accept 18 of 100 = coverage 0.18.
\item \textbf{Precision}: of the accepted ones, what fraction are
\emph{actually correct}. Example: 182 correct of 184 accepted = 0.989.
\end{itemize}
\keybox{High precision on accepted items is the whole point: the gate
should only wave through receipts it is almost never wrong about.}
\end{frame}

\begin{frame}{Sample size $n$ and why it rules everything}
\plain{$n$ is just how many examples you measured.}
\begin{itemize}
\item 15 out of 15 correct sounds perfect. But with only $n=15$ you
cannot tell ``truly excellent'' from ``lucky''.
\item More $n$ = a sharper estimate. Tiny $n$ = a blurry one, even at 100\%.
\end{itemize}
\keybox{Small $n$ does not make a result wrong. It makes it
\emph{uninformative}: consistent with both ``great'' and ``mediocre''.}
\end{frame}

\begin{frame}{Confidence interval and ``lower bound''}
\plain{A confidence interval is the honest range the true value could be in.}
\begin{itemize}
\item 15/15 at $n=15$: best guess 1.000, but the 95\% \emph{lower bound}
is only about \textbf{0.796} (Wilson). The truth could be as low as $\sim$0.80.
\item 113/114 at $n=472$: lower bound about \textbf{0.95}. Much tighter,
because $n$ is larger.
\end{itemize}
\keybox{We always quote the \emph{lower bound}, the pessimistic honest
floor, not the flattering point estimate.}
\end{frame}

\begin{frame}{Statistical power and ``underpowered''}
\plain{Power = the chance your test can \emph{detect} a real effect that
is truly there.}
\begin{itemize}
\item If $n$ is too small, even a real effect fails the test: not
because it is absent, but because the test had no resolution.
\item ``\textbf{Underpowered}'' = ``$n$ too small to conclude'', NOT
``the method failed''. These are very different sentences.
\end{itemize}
\keybox{An underpowered cell is honest evidence of \emph{not knowing
yet}, not evidence \emph{against}.}
\end{frame}

\begin{frame}{TOST: a fair pass/fail at a threshold}
\plain{TOST asks one crisp question: is precision provably above a bar
(say 0.95)?}
\begin{itemize}
\item It returns a $p$-value. Small $p$ (say $<0.05$) = ``yes, provably
above 0.95''. Large $p$ = ``cannot show it at this $n$''.
\item It is deliberately strict: it refuses to pass on small samples.
\end{itemize}
\keybox{TOST is the referee. We report exactly which corpora it passes
and which it does not. No spin.}
\end{frame}

\begin{frame}{AUROC and ``baseline''}
\plain{AUROC = how well a single score separates two groups. 1.0 perfect,
0.5 useless coin-flip.}
\begin{itemize}
\item A \textbf{baseline} is a simple existing signal you must beat or at
least match (here: softmax confidence, entropy, sequence length, etc.).
\item If your score's AUROC is \emph{below} every baseline, it is worse
at that specific job than things people already have.
\end{itemize}
\keybox{Whether that is fatal depends entirely on \emph{what job} you
claim to be doing. Hold that thought.}
\end{frame}

\begin{frame}{Distribution shift (Paper 2's setting)}
\plain{A model trained on one kind of receipt is run on a different
kind. That change is ``distribution shift''.}
\begin{itemize}
\item In-distribution: familiar receipts (CORD).
\item Shifted: different receipts (SROIE).
\item Question: can we \emph{detect} that we are now off-distribution,
using only the model's own outputs, no labels?
\end{itemize}
\keybox{Paper 2 is about \emph{detecting} shift, not about fixing or
triaging individual predictions. Hold that too.}
\end{frame}

% =====================================================================
\section{Part 2: Caveat 1 (Paper 1) is per-corpus underpowered}
% =====================================================================
\begin{frame}{Paper 1 in one line}
\bigq{An arithmetic check (sigma) and the model's softmax confidence,
\emph{intersected}, accept a receipt only when \textbf{both} agree.}
\vspace{0.8em}
\plain{Where both agree, precision is very high. The question is: how
strong is that evidence, corpus by corpus?}
\end{frame}

\begin{frame}{The real numbers (intersection precision)}
\begin{itemize}
\item \textbf{Pooled}: $182/184 = \textbf{0.989}$, 95\% lower bound
\textbf{0.961} (corpus-stratified bootstrap CI about $[0.973,1.000]$).
\item \textbf{WildReceipt} ($n_{\text{corpus}}=472$): $113/114=0.991$,
lower bound $\approx 0.95$.
\item \textbf{CORD} ($n=200$): $54/55=0.982$.
\item \textbf{SROIE} ($n=347$): $15/15=1.000$, but only \textbf{15}
accepted.
\end{itemize}
\keybox{Looks strong. Now look at \emph{where the evidence lives}.}
\end{frame}

\begin{frame}{The concentration problem (basic)}
\plain{Of the 184 pooled intersection accepts, \textbf{114} come from
WildReceipt alone.}
\vspace{0.6em}
\begin{itemize}
\item That is about \textbf{62\%} of the headline mass in one corpus.
\item SROIE contributes only \textbf{15} of 184 ($\approx 8\%$).
\end{itemize}
\keybox{The pooled 0.989 is real, but it is \emph{carried by
WildReceipt}. A reviewer will, correctly, recompute this.}
\end{frame}

\begin{frame}{The small-$n$ problem (intermediate)}
SROIE is $15/15$. Its honest floors at $n=15$:
\begin{itemize}
\item Wilson 95\% lower bound: \textbf{0.796}
\item Clopper--Pearson exact: \textbf{0.819}
\item Jeffreys: \textbf{0.882}
\end{itemize}
\vspace{0.4em}
\keybox{Even a \emph{perfect} 15/15 only guarantees ``about 0.80 or
better''. That is a \emph{consistency} statement, not strong evidence.}
\end{frame}

\begin{frame}{The referee's verdict (advanced): per-corpus TOST vs 0.95}
One-sided exact-binomial TOST, threshold 0.95, $\alpha=0.05$:
\begin{itemize}
\item \textbf{WildReceipt}: $p=\textbf{0.020}$ \quad
\textcolor{ok}{passes}
\item \textbf{CORD}: $p=0.232$ \quad \textcolor{bad}{does not pass}
\item \textbf{SROIE}: $p=0.463$ \quad \textcolor{bad}{does not pass}
\end{itemize}
\vspace{0.3em}
Leave-one-corpus-out minimum lower bound: \textbf{0.923} (still strong
pooled). Within-corpus error-decorrelation still holds everywhere
(Stouffer $p\approx 7\times 10^{-11}$, Fisher $p\approx 3\times 10^{-10}$).
\keybox{Honest reading: the per-corpus claim is individually proven only
on WildReceipt; CORD and SROIE are \emph{consistent but underpowered}.}
\end{frame}

\begin{frame}{Why editing cannot fix Caveat 1}
\plain{To make CORD/SROIE individually pass, you do not need better
words. You need more accepted receipts.}
\begin{itemize}
\item Required $n$ for 80\% power: CORD $\approx$\textbf{242},
SROIE $\approx$\textbf{71}, WildReceipt already there ($\approx$161).
\item That means: add corpora with enough intersection mass (DocILE,
mychen76, ...) \emph{and} GPU fine-tune the models on them.
\item Not possible in this editing interface. And even done, it raises
\emph{power}; it does not guarantee a reviewer \emph{values} it.
\end{itemize}
\keybox{So we do the only honest thing: report the floors, the
leave-one-out, and the exact TOST pass/fail, plainly.}
\end{frame}

% =====================================================================
\section{Part 3: Caveat 2 (Paper 2) is below every baseline on raw AUROC}
% =====================================================================
\begin{frame}{Paper 2 in one line}
\bigq{Under distribution shift the beam-margin's \emph{variance}
collapses by about \textbf{169}$\times$ (log$_2 C = 7.40$), a change the
classic KS test barely sees (KS $=0.735$).}
\vspace{0.6em}
\plain{That variance-collapse \emph{law} is the contribution. Now the
uncomfortable table.}
\end{frame}

\begin{frame}{The uncomfortable table (real AUROC numbers)}
How well does each signal, used \emph{per receipt}, detect shift?
\begin{itemize}
\item Beam-margin (ours): \textbf{0.875}
\item Sequence length: 0.986 \quad softmax confidence $c_{\text{seq}}$: 0.985
\item Token entropy: 0.984 \quad MSP: 0.980 \quad Mahalanobis: 0.955
\item (Free energy: 0.639)
\end{itemize}
\vspace{0.3em}
\keybox{Used as a raw per-receipt detector, the margin is \emph{below
every serious baseline}. This is stated in the paper, not hidden.}
\end{frame}

\begin{frame}{Why this is intrinsic, not a bug (intermediate)}
\plain{The margin was never a per-receipt score. It is a
\emph{population} statistic: a property of a \emph{batch}'s variance.}
\begin{itemize}
\item Asking ``what is its per-receipt AUROC'' is asking a question it
is not built to answer.
\item Its value is elsewhere: a \emph{characterised law} (169$\times$
variance collapse) with a \emph{mechanism} (entropy-floor saturation),
label-free, computed from outputs the model already produced.
\end{itemize}
\keybox{The KS test, the standard tool, sees only 0.735 and misses the
169$\times$ collapse entirely. That gap \emph{is} the finding.}
\end{frame}

\begin{frame}{Why editing cannot fix Caveat 2 (advanced)}
\begin{itemize}
\item To make the margin beat the baselines on raw AUROC, the signal
would have to be a strong per-example score. \textbf{It is not}; that is
its nature.
\item The only way to ``remove'' this caveat by editing is to
\emph{claim} it beats baselines. That is exactly the dishonesty this
whole effort removed.
\item So we keep it: state the AUROC table straight, and locate the
contribution where it actually is (the law + mechanism + inference-only,
label-free deployment monitor).
\end{itemize}
\keybox{A true sentence that limits your claim is not a defect. A false
sentence that inflates it is.}
\end{frame}

% =====================================================================
\section{Part 4: Caveat 3 (Paper 2): no triage value, single shift pair}
% =====================================================================
\begin{frame}{The honest negative, stated by the paper itself}
\plain{We ran the controlled test: is the per-receipt margin useful for
deciding \emph{which single prediction} to trust?}
\begin{itemize}
\item Risk (AURC, lower better): margin \textbf{0.513} vs softmax
$c_{\text{seq}}$ \textbf{0.516}. Effectively \emph{no advantage}.
\item On the shifted set, error rate 66\%; margin recall 23\% vs
$c_{\text{seq}}$ 22\%.
\end{itemize}
\keybox{Conclusion in the paper: this is a shift \emph{detector} with a
mechanism, \textbf{not} an operational per-receipt triage tool. We do
not claim triage value.}
\end{frame}

\begin{frame}{The single-shift-pair point (intermediate)}
The headline rests on one natural pair: CORD $\rightarrow$ SROIE. We add
three independent corroborations:
\begin{itemize}
\item Within-CORD natural complexity shift: log$_2 C = 3.63$
\item A third real receipt/invoice corpus: log$_2 C = 7.53$
\item Held-out, leakage-free re-measurement: log$_2 C = 7.40$
\item Model axis: \textbf{2 of 3} sibling checkpoints fire
\end{itemize}
\keybox{Stronger than one pair alone, but still not a broad multi-shift
validation. We say exactly that.}
\end{frame}

\begin{frame}{Why editing cannot fix Caveat 3 (advanced)}
\begin{itemize}
\item ``Remove the no-triage-value caveat'' would require the margin to
\emph{actually be} an operationally useful triage score. The controlled
test says it is not (AURC 0.513 vs 0.516).
\item ``Remove the single-pair caveat'' would require validation on many
\emph{natural} shift pairs: new datasets plus new GPU model runs. That
is a different, larger paper, not an edit.
\item Deleting either sentence without doing that work would be a
false claim.
\end{itemize}
\keybox{The integrity of Paper 2 \emph{is} that it argues its own
limits. Removing them removes the honesty, not the limitation.}
\end{frame}

% =====================================================================
\section{Part 5: the significance numbers behind Paper 2}
% =====================================================================
\begin{frame}{For completeness: the recomputed leakage-free stats}
After removing training-data contamination (headline now $n_{\text{in}}=200$
held-out vs SROIE 347):
\begin{itemize}
\item log$_2 C = 7.40$ (variance ratio $\approx 169\times$); KS only 0.735.
\item Levene/Brown--Forsythe $W=206.45$; Bartlett $T=1229.87$; $F=169$.
\item Bootstrap 95\% CI on log$_2 C$: $[6.64,\,8.24]$ (excludes 0).
\item Permutation null max $|\log_2 C|=1.99$; the observed 7.40 sits
$\approx 3.7\times$ outside it; $p<10^{-4}$.
\end{itemize}
\keybox{The effect is real and robust on the \emph{correct} data. The
caveats above are about \emph{scope and per-example use}, not about
whether the law holds. It holds.}
\end{frame}

% =====================================================================
\section{Part 6: the meta-point}
% =====================================================================
\begin{frame}{What we fixed vs what we kept}
\begin{itemize}
\item \textcolor{ok}{Fixed (copy-edit defects):} a contaminated
evaluation, a figure with stale numbers, an internally contradictory
paragraph, a fabricated citation, em-dashes, column overflows.
\item \textcolor{bad}{Kept and stated (empirical truths):} per-corpus
underpowered (Paper 1); below baselines on raw AUROC (Paper 2); no
per-receipt triage value and a single main shift pair (Paper 2).
\end{itemize}
\keybox{The first list is now empty. The second list is, deliberately,
written into the papers in plain language.}
\end{frame}

\begin{frame}{Why this is the strong position, not the weak one}
\begin{itemize}
\item A reviewer who finds a hidden weakness rejects you and distrusts
the rest. A reviewer who sees you state it first trusts the rest.
\item Every number here is the real measured value, including the
unflattering ones.
\item No paper at any venue is a guaranteed accept. Stating limits
precisely is the most defensible posture for an applied venue.
\end{itemize}
\keybox{Honest, reproducible, and exactly scoped: no more, and no less,
than the data supports.}
\end{frame}

\begin{frame}{One-slide summary}
\begin{itemize}
\item \textbf{Caveat 1} (Paper 1): pooled 0.989 is carried by
WildReceipt; CORD/SROIE pass only as ``consistent, underpowered''
(TOST $p$: 0.020 / 0.232 / 0.463). Fix needs more data + GPU, not edits.
\item \textbf{Caveat 2} (Paper 2): per-receipt AUROC 0.875, below every
baseline, because the signal is a population law, not a per-example
score. Cannot be edited away without lying.
\item \textbf{Caveat 3} (Paper 2): no triage value (AURC 0.513 vs
0.516), one main shift pair plus 3 corroborations. Removing it needs new
science, not wording.
\end{itemize}
\keybox{They are facts. We state them. That is the contribution's
credibility.}
\end{frame}

\end{document}
